% p30-f-natural-completes-the-following-diagram.tex
%
% Source:   Carl A. Gunter and Dana S. Scott, "Semantic Domains".
%           ScottDomains/papers/Gunter Scott 1990.pdf
% Physical page: 30          Printed page: 29
% Section:  5.3 Universal and closure properties.
% Unnamed in the paper (no figure number). Introduced by:
%   "As was the case with the smash product and lift operators, a diagram like
%    the one in Theorem 12 gives rise to another important operation on
%    functions. If f : D -> E is a continuous function, then there is a unique
%    homomorphism f-natural which completes the following diagram:"
%
% Content: a commutative square: the singleton map {|.|} is natural in D, and
% f-natural is the lifted action of the convex powerdomain on morphisms. The
% paper adds: "Namely, one defines f-natural = ext({|.|} o f)."
% Vertices 4, arrows 4 (3 solid, 1 dashed).
%
% Geometry measured from a 300 dpi render of the region and divided by
% 118.11 px/cm.
%
\documentclass[border=6pt]{standalone}
\usepackage{tikz}
% the paper's singleton bracket {|.|}; the negative kerns must stay small
% enough that the vertical bars are not swallowed by the braces
\newcommand{\sglt}{\{\mkern-3.5mu|\mkern1mu\cdot\mkern1mu|\mkern-3.5mu\}}
\begin{document}
\begin{tikzpicture}[
  x=1cm, y=1cm,
  every node/.style={inner sep=1pt},
  open/.style={circle, draw, fill=white, minimum size=4pt, inner sep=0pt},
  solid/.style={circle, draw, fill=black, minimum size=5pt, inner sep=0pt},
  edge/.style={draw, line width=0.4pt},
  % added for the commutative diagrams: object nodes, and the paper's two
  % arrow kinds -- solid for a named map, dashed for the unique completing map
  obj/.style={inner sep=3pt},
  arrow/.style={draw, line width=0.4pt, ->},
  darrow/.style={draw, line width=0.4pt, ->, dash pattern=on 4pt off 3pt}]

  % vertices
  \node[obj] (d)   at (0.00,2.07) {$D$};
  \node[obj] (e)   at (3.58,2.07) {$E$};
  \node[obj] (dna) at (0.00,0.00) {$D^{\natural}$};
  \node[obj] (ena) at (3.58,0.00) {$E^{\natural}$};

  % arrows
  \draw[arrow]  (d)   -- (e);
  \draw[arrow]  (d)   -- (dna);
  \draw[arrow]  (e)   -- (ena);
  \draw[darrow] (dna) -- (ena);

  % labels printed inside the picture
  \node[anchor=south] at (1.79,2.19) {$f$};
  \node[anchor=south] at (1.79,0.19) {$f^{\natural}$};
  \node[anchor=east]  at (-0.25,1.02) {$\sglt$};
  \node[anchor=west]  at ( 3.83,1.02) {$\sglt$};
\end{tikzpicture}
\end{document}
